\documentclass[../mainReport.tex]{subfiles} % 继承主文件设置
\begin{document}
\section{Report1}\label{sec: report1}
\subsection{rp++}\label{sec: rp++}
\subsubsection{形式化定义}\label{sec: rp++形式化定义}
给定字符串$S = s_1s_2...s_n$，其中$n \leq 2 \times 10^6$，定义模式：
$$
P = \text{rp++} = [r][p][+][+]
$$
要求统计$P$在$S$中的所有精确匹配次数

% \subsubsection{完整代码}\label{sec: 完整代码}
% 这题水水的，唯一需要注意的是\textcolor{blue}{记得更新变量}\verb|i += t.length();|，不然程序会死循环。

% \begin{lstlisting}[caption={题解1}]
% #include<bits/stdc++.h>
% using namespace std;

% int main(){
% 	int cnt = 0;
% 	string s, t = "rp++";
% 	int len = t.length();
% 	cin >> s;
% 	string::size_type i = 0;
% 	while((i = s.find(t, i)) != string::npos){
% 		cnt++;
% 		i += len;
% 	}
% 	printf("%d", cnt);
%     return 0;
% }

% \end{lstlisting}

\subsection{数组游戏}
\subsubsection{形式化定义}
给定初始空数组$ A_0 = \emptyset $，经过$q$次操作后得到数组序列$\{A_t\}_{t=1}^q$，定义第$t$次操作后的\textbf{炫酷值}为：
$$
F(A_t) = \sum_{i=1}^{|A_t|} A_t[i] \cdot i
$$

操作类型集合$\mathcal{O} = \{1,2,3\}$定义如下：
\begin{itemize}
\item $\text{op}=1$：循环右移，$A_t = \text{rotate}(A_{t-1}) = [a_n,a_1,...,a_{n-1}]$
\item $\text{op}=2$：反转数组，$A_t = \text{reverse}(A_{t-1}) = [a_n,...,a_1]$
\item $\text{op}=3\ (k)$：追加元素，$A_t = A_{t-1} \circ [k]$
\end{itemize}

对于复杂度：

\begin{itemize}
\item 操作次数：$q \leq 2 \times 10^5$
\item 追加数值：$k \leq 10^6$
\item 时间限制：通常为1-2秒（需$O(1)$或$O(\log n)$每次操作）
\end{itemize}

\subsubsection{解题思路}
本解法采用\textbf{双端队列+标记法}高效维护动态数组操作，核心思想如下：

\begin{enumerate}
\item \textbf{状态维护}：
\begin{itemize}
\item 使用双端队列$dq$存储实际数组元素
\item 标记位$reversed$记录当前是否处于反转状态
\item 维护三个关键变量：
\begin{align*}
S &= \sum_{i=1}^n (a_i \times i) \quad \text{(当前炫酷值)} \\
T &= \sum_{i=1}^n a_i \quad \text{(元素总和)} \\
n &= |dq| \quad \text{(数组长度)}
\end{align*}
\end{itemize}

\item \textbf{操作处理}(经数学推导可证明)：
\begin{itemize}
\item \textcolor{blue}{循环右移（op=1）}：

因为：

\begin{equation*}
    S'=a_n \cdot 1 + \sum_{i=1}^{n-1}a_i\cdot(i+1)=S+\sum_{i=1}^{n}a_i-n\cdot a_n\\
\end{equation*}

所以有：

\begin{equation*}
S \leftarrow S + \underbrace{T - n \times a_{\text{移动元素}}}_{\text{增量调整}}
\end{equation*}
根据$reversed$状态决定移动首部或尾部元素

\item \textcolor{blue}{数组反转（op=2）}：

因为：

\begin{equation*}
    S'=\sum_{i=1}^na_{n-i+1}\cdot i=\sum_{j=1}^na_j\cdot(n-j+1)=(n+1)T-S\\
\end{equation*}

所以有：

\begin{equation*}
S \leftarrow (n+1) \times T - S
\end{equation*}
通过数学变换避免实际反转操作

\item \textcolor{blue}{追加元素（op=3）}：

显然有：

\begin{equation*}
\begin{cases}
S \leftarrow S + k \times (n+1) \\
T \leftarrow T + k \\
n \leftarrow n + 1
\end{cases}
\end{equation*}
根据$reversed$状态决定插入方向
\end{itemize}
\end{enumerate}

\subsubsection{复杂度分析}
\begin{table}[h]
\centering
\begin{tabular}{|l|c|}
\hline
操作类型 & 时间复杂度 \\
\hline
循环移位 & $O(1)$ \\
数组反转 & $O(1)$ \\
追加元素 & $O(1)$ \\
查询炫酷值 & $O(1)$ \\
\hline
\end{tabular}
\end{table}


% \subsubsection{完整代码}

% \begin{lstlisting}[caption={题解2}]
% #include <bits/stdc++.h>
% using namespace std;

% int main() {
%     int q;
%     cin >> q;
%     deque<int> dq;
%     long long S = 0, T = 0, n = 0;
%     bool reversed = false; // 标记是否反转

%     while (q--) {
%         int op;
%         cin >> op;
%         if (op == 1) { // 循环右移1位
%             if (!reversed) {
%                 int last = dq.back();
%                 dq.pop_back();
%                 dq.push_front(last);
%                 S += T - n * last;
%             } else {
%                 int first = dq.front();
%                 dq.pop_front();
%                 dq.push_back(first);
%                 S += T - n * first;
%             }
%         } else if (op == 2) { // 反转数组
%             reversed = !reversed;
%             S = (n + 1) * T - S;
%         } else if (op == 3) { // 追加元素k
%             int k;
%             cin >> k;
%             if (!reversed) {
%                 dq.push_back(k);
%             } else {
%                 dq.push_front(k);
%             }
%             n++;
%             T += k;
%             S += k * n;
%         }
%         printf("%lld\n", S);
%     }
%     return 0;
% }

% \end{lstlisting}

\subsection{你喜欢Inf、并查集和算法实验吗}

\subsubsection{形式化定义}
给定社交关系图 $G=(V, E_f \cup E_h)$，其中：
\begin{itemize}
    \item $V$ 表示所有熟人节点（$|V|=n$）
    \item $E_f$ 为友谊边集（无向边）
    \item $E_h$ 为敌对边集（无向边）
    \item 满足 $E_f \cap E_h = \emptyset$
\end{itemize}

一个合法派对 $P \subseteq V$ 需满足：
\begin{enumerate}
    \item \textbf{朋友闭包性}：$\forall v \in P$，其所有朋友必须被邀请
    $$ N_f(v) \subseteq P \quad \text{其中} \quad N_f(v) = \{ u \mid (v,u) \in E_f \} $$
    
    \item \textbf{无敌对约束}：派对内不存在互相敌对的两人
    \begin{equation*}
        \not \exists (u,v) \in E_h \quad \text{s.t.} \quad u,v \in P 
    \end{equation*}
    
    \item \textbf{连通性}：派对成员在友谊图中连通
    $$ \forall u,v \in P, \exists \text{路径} \ u \leftrightarrow v \ \text{仅通过} \ E_f \ \text{中的边} $$
\end{enumerate}

\subsubsection{输入输出要求}
\begin{itemize}
    \item 第一行：$n$（熟人数量，$2 \leq n \leq 2000$）
    \item 第二行：$k$（友谊边数，$0 \leq k \leq \min(10^5, \frac{n(n-1)}{2})$）
    \item 随后$k$行：每行$u_i\ v_i$表示一对朋友
    \item 下一行：$m$（敌对边数，$0 \leq m \leq \min(10^5, \frac{n(n-1)}{2})$）
    \item 随后$m$行：每行$u_j\ v_j$表示一对敌对关系
\end{itemize}

寻找满足条件的最大派对人数，若不存在合法解则输出0。

\begin{verbatim}
输入：
9
8
1 2
1 3
2 3
4 5
6 7
7 8
8 9
9 6
2
1 6
7 9
\end{verbatim}

\begin{itemize}
    \item 有效解：$\{1,2,3\}$（大小3）和$\{4,5\}$（大小2）
    \item 无效解：$\{6,7,8,9\}$（违反敌对约束），$\{1,2,3,4,5\}$（违反连通性）
    \item 输出：3
\end{itemize}

\subsubsection{并查集应用}
我们注意到输入数据的特点，\textcolor{blue}{每对人最多在输入中被提到一次，两个人不能同时是朋友和彼此不喜欢}，所以我们可以考虑用\textcolor{blue}{并查集}来管理不同的朋友链。

而在\verb|unite()|的过程中，需要注意维护当前朋友链的大小\verb|Size[n]|，结合\textcolor{blue}{按秩合并}的思想，我们可以用这个数组来指导不同的朋友链之间的合并。

% \begin{lstlisting}[caption={题解3：并查集部分}]
% int find(int x)
% {
%     if (parent[x] != x)
%     {
%         parent[x] = find(parent[x]);
%     }
%     return parent[x];
% }

% void unite(int x, int y)
% {
%     int rootX = find(x);
%     int rootY = find(y);
%     if (rootX != rootY)
%     {
%         if (Size[rootX] < Size[rootY])
%         {
%             swap(rootX, rootY);
%         }
%         parent[rootY] = rootX;
%         Size[rootX] += Size[rootY];
%     }
% }

% \end{lstlisting}

\subsubsection{敌对关系检查}
接下来我们进行\textcolor{blue}{敌对关系检查}，只需要逐一验证每个朋友链中是否存在敌对关系即可，如果存在，那么将这个朋友链标记为\verb|false|即可。

% \begin{lstlisting}[caption={题解3：敌对关系检查部分}]
% scanf("%d", &m);
% for (int i = 1; i <= m; i++)
% {
%     int u, v;
%     cin >> u >> v;
%     dislike[u].insert(v);
%     dislike[v].insert(u);
% }

% for (int i = 1; i <= n; i++)
% {
%     int root = find(i);
%     if (!valid[root]) // 已经被否定的朋友链不再遍历
%         continue;
%     for (int j = 1; j <= n; j++)
%     {
%         if (find(j) == root) // 保证 j 与 i 在同一个朋友链中
%         {
%             for (int d : dislike[j])
%             {
%                 if (find(d) == root) // 如果 j 的敌对关系也在这个朋友链中，则否定该朋友链
%                 {
%                     valid[root] = false;
%                     break;
%                 }
%             }
%             if (!valid[root]) // 已经被否定的朋友链不再遍历 
%                 break;
%         }
%     }
% }
% \end{lstlisting}

% \subsubsection{完整代码}

% \begin{lstlisting}
% #include <bits/stdc++.h>
% using namespace std;

% const int MAXN = 2005;
% int parent[MAXN];
% int Size[MAXN]; // 不与size重复
% bool valid[MAXN];
% int n, k, m;
% unordered_set<int> dislike[MAXN];
% int find(int x)
% {
%     if (parent[x] != x)
%     {
%         parent[x] = find(parent[x]);
%     }
%     return parent[x];
% }

% void unite(int x, int y)
% {
%     int rootX = find(x);
%     int rootY = find(y);
%     if (rootX != rootY)
%     {
%         if (Size[rootX] < Size[rootY])
%         {
%             swap(rootX, rootY);
%         }
%         parent[rootY] = rootX;
%         Size[rootX] += Size[rootY];
%     }
% }

% int main()
% {
%    // freopen("../in.txt", "r", stdin);
%     scanf("%d%d", &n, &k);

%     for (int i = 1; i <= n; i++)
%     {
%         parent[i] = i;
%         Size[i] = 1;
%         valid[i] = true;
%     }

%     for (int i = 1; i <= k; i++)
%     {
%         int u, v;
%         scanf("%d %d", &u, &v);
%         unite(u, v);
%     }

%     scanf("%d", &m);
%     for (int i = 1; i <= m; i++)
%     {
%         int u, v;
%         cin >> u >> v;
%         dislike[u].insert(v);
%         dislike[v].insert(u);
%     }

%     for (int i = 1; i <= n; i++)
%     {
%         int root = find(i);
%         if (!valid[root])
%             continue;
%         for (int j = 1; j <= n; j++)
%         {
%             if (find(j) == root)
%             {
%                 for (int d : dislike[j])
%                 {
%                     if (find(d) == root)
%                     {
%                         valid[root] = false;
%                         break;
%                     }
%                 }
%                 if (!valid[root]) 
%                     break;
%             }
%         }
%     }

%     int mx = 0;
%     for (int i = 1; i <= n; i++)
%     {
%         int root = find(i);
%         if (valid[root] && Size[root] > mx)
%         {
%             mx = Size[root];
%         }
%     }

%     printf("%d", mx);
%     return 0;
% }
% \end{lstlisting}

\subsection{Sherway的有向图}
\subsubsection{形式化定义}

给定有向完全图$G=(V,E)$，其中：
\begin{itemize}
\item 顶点集$V=\{1,2,\ldots,n\}$
\item 边集$E$由邻接矩阵$a$定义，$a_{i,j}$表示边$(i,j)$的权重
\item 初始所有顶点处于关闭状态$C=V$，开放顶点集$O=\emptyset$
\end{itemize}

定义状态转移函数：
$$
\begin{cases}
O \leftarrow O \cup \{x\} & \text{操作1 x} \\
O \leftarrow O \setminus \{x\} & \text{操作3（撤销最近操作1）}
\end{cases}
$$

对于操作2 x y，定义：
$$
\delta(x,y) = \min_{\substack{p \in P(x,y) \\ p \cap C = \emptyset}} \sum_{(u,v)\in p} a_{u,v}
$$
其中$P(x,y)$为$x$到$y$的所有路径集合。

\begin{itemize}
\item 邻接矩阵满足$a_{i,i}=0,\ \forall i\in V$
\item 操作约束：
  \begin{itemize}
  \item 操作1前$x \in C$
  \item 操作2前$x,y \in O$
  \item 操作3前存在未撤销的操作1
  \end{itemize}
\end{itemize}

\begin{itemize}
\item 空间复杂度：$O(n^2)$存储邻接矩阵
\item 时间复杂度：
  \begin{itemize}
  \item 操作1/3：$O(1)$
  \item 操作2：$O(n^2)$（Floyd-Warshall动态更新）
  \end{itemize}
\end{itemize}


\subsubsection{floyd改进}
题目的数据范围限定了用经典的复杂度为$O(n^3)$的算法，会超时（只过两个点），我们要想到，大部分路线是没必要重新计算的，新引入一个点带来的影响的路径，\textcolor{blue}{只可能是以这个点为中间结点的路径}。

% \begin{lstlisting}[caption={题解4：floyd改进部分}]
% // 原本的floyd
% void floyd()
% {
%     for (int k = 1; k <= n; ++k)
%     {
%         if (!open[k])
%             continue;
%         for (int i = 1; i <= n; ++i)
%         {
%             if (!open[i])
%                 continue;
%             for (int j = 1; j <= n; ++j)
%             {
%                 if (!open[j])
%                     continue;
%                 if (cur[i][k] + cur[k][j] < cur[i][j])
%                 {
%                     cur[i][j] = cur[i][k] + cur[k][j];
%                 }
%             }
%         }
%     }
% }

% // 改进后的floyd
% void floyd2(int x)
% {
%     for (int i = 1; i <= n; ++i)
%     {
%         for (int j = 1; j <= n; ++j)
%         {
%             if (cur[i][j] > cur[i][x] + cur[x][j])
%             {
%                 cur[i][j] = cur[i][x] + cur[x][j];
%             }
%         }
%     }
% }
% \end{lstlisting}

\subsubsection{操作3优化}

因为操作3不是随机取消一个点，而是撤销上一个操作1，因此我们可以\textcolor{blue}{用一个栈来维护被操作1更新过的图}，这样我们就可以不再计算一次floyd算法，提高效率。

当然这里如果不采用\verb|cur|数组来维护，全局只用一个栈，栈顶表示当前图也是可以的，这样在空间复杂度不变的情况下，避免了上面操作3的拷贝操作带来的时间损耗，可以提高运行效率。

% \begin{lstlisting}[caption={题解4： 操作3优化部分}]
% if (op == 1)
% {
%     int x;
%     scanf("%d", &x);
%     open[x] = true;
%     op1s.push(x);

%     st_mp.push(vector<vector<int>>(n + 1, vector<int>(n + 1)));
%     for (int i = 1; i <= n; ++i)
%         for (int j = 1; j <= n; ++j)
%             st_mp.top()[i][j] = cur[i][j];
%     floyd2(x);
% }
% else if (op == 2)
% {
%     int x, y;
%     scanf("%d %d", &x, &y);
%     printf("%d\n", cur[x][y]);
% }
% else if (op == 3)
% {
%     auto &backup = st_mp.top();
%     for (int i = 1; i <= n; ++i)
%         for (int j = 1; j <= n; ++j)
%             cur[i][j] = backup[i][j];
%     st_mp.pop();
%     open[op1s.top()] = false;
%     op1s.pop();
% }
% \end{lstlisting}

\subsubsection{复杂度分析}

\begin{table}[h]
\centering
\begin{tabular}{|l|c|}
\hline
操作类型 & 时间复杂度 \\
\hline
打开结点 & $O(n^2)$ \\
查询最短路径 & $O(1)$ \\
撤销操作1 & $O(n^2)$(可优化为$O(1)$) \\
\hline
\end{tabular}
\end{table}

所以本题的总体时间复杂度为$O(t_s\cdot n^2)$，空间复杂度为$O(t_1\cdot n^2)$；可以优化为时间复杂度为$O(t_1\cdot n^2)$，空间复杂度为$O(t_1\cdot n^2)$。($t_s$是操作1,3的总数量，$t_1$是操作1的数量)
% \subsubsection{完整代码}

% \begin{lstlisting}
% #include <bits/stdc++.h>
% #define SHI_YAN_KE_JIE_SHU_LE 0
% using namespace std;

% const int INF = 1e9;
% const int MAXN = 305;

% int n, q;
% int cur[MAXN][MAXN];
% bool open[MAXN];
% stack<int> op1s; 
% stack<vector<vector<int>>> st_mp;


% void floyd2(int x)
% {
% 	for (int i = 1; i <= n; ++i)
% 	{
% 		for (int j = 1; j <= n; ++j)
% 		{
% 			if (cur[i][j] > cur[i][x] + cur[x][j])
% 			{
% 				cur[i][j] = cur[i][x] + cur[x][j];
% 			}
% 		}
% 	}
% }

% int main()
% {
% 	// freopen("in.txt", "r", stdin);
% 	scanf("%d %d", &n, &q);
% 	for (int i = 1; i <= n; ++i)
% 	{
% 		for (int j = 1; j <= n; ++j)
% 		{
% 			scanf("%d", &cur[i][j]);
% 		}
% 	}
% 	memset(open, false, sizeof(open));
% 	while (q--)
% 	{
% 		int op;
% 		scanf("%d", &op);
% 		if (op == 1)
% 		{
% 			int x;
% 			scanf("%d", &x);
% 			open[x] = true;
% 			op1s.push(x);

% 			st_mp.push(vector<vector<int>>(n + 1, vector<int>(n + 1)));
% 			for (int i = 1; i <= n; ++i)
% 				for (int j = 1; j <= n; ++j)
% 					st_mp.top()[i][j] = cur[i][j];
% 			floyd2(x);
% 		}
% 		else if (op == 2)
% 		{
% 			int x, y;
% 			scanf("%d %d", &x, &y);
% 			printf("%d\n", cur[x][y]);
% 		}
% 		else if (op == 3)
% 		{
% 			auto &backup = st_mp.top();
% 			for (int i = 1; i <= n; ++i)
% 				for (int j = 1; j <= n; ++j)
% 					cur[i][j] = backup[i][j];
% 			st_mp.pop();
% 			open[op1s.top()] = false;
% 			op1s.pop();
% 		}
% 	}
% 	return SHI_YAN_KE_JIE_SHU_LE;
% }
% \end{lstlisting}


\end{document}